\documentclass{beamer}
\usetheme{Madrid}
\usecolortheme{default}

\title{Proyecto 1 - OLTP - Colecciones}
\subtitle{Bases de Datos a Gran Escala}
\author{Equipo 13: \\ Andres Lopez Gonzalez \\ Carlos Fabian Esteva Gallegos}
\date{}

\begin{document}

%---------------------------
\begin{frame}
  \titlepage
\end{frame}

%---------------------------
\begin{frame}{Descripción del Proyecto}
  \begin{itemize}
    \item Encuentro mundial de coleccionistas de tarjetas deportivas.
    \item Se diseña un sistema OLTP para gestionar deportes, equipos, jugadores y tarjetas.
    \item Información clave:
    \begin{itemize}
      \item Deporte: nombre, descripción.
      \item Equipo: nombre, ciudad, fechas de partidos.
      \item Jugador: nombre, género, puesto, fechas, lugar de nacimiento.
      \item Tarjeta: número, estado, fecha venta, medio adquisición.
    \end{itemize}
  \end{itemize}
\end{frame}

%---------------------------
\begin{frame}{Modelo y Reglas}
  \begin{itemize}
    \item Se aplicaron \textbf{constraints} para integridad:
    \begin{itemize}
      \item Validación de nombres, géneros y estados.
      \item Reglas de consistencia en fechas (ej. nacimiento $<$ pertenencia equipo).
      \item Restricciones de adquisición y estado de tarjetas.
    \end{itemize}
    \item Procedimientos almacenados:
    \begin{itemize}
      \item \texttt{sp\_inserta}: inserción de deportes.
      \item \texttt{sp\_borra}: eliminación segura de deportes (con validación de dependencias).
    \end{itemize}
  \end{itemize}
\end{frame}

%---------------------------
\begin{frame}{Consultas y Resultados}
  \begin{itemize}
    \item \texttt{sp\_consultas} implementa 5 consultas parametrizadas:
    \begin{enumerate}
      \item Estado de tarjetas por equipo.
      \item Medio de adquisición por jugador.
      \item Relación deporte-equipo-tarjeta.
      \item Equipos según estado de colección.
      \item Jugadores y puesto dentro de un equipo.
    \end{enumerate}
    \item Resultados exportados a archivos de texto (\texttt{.txt}) para su análisis.
  \end{itemize}
\end{frame}

\end{document}

